\documentclass[11pt]{amsart}
\usepackage{amsmath,amssymb,amsthm}
\usepackage[margin=1.05in]{geometry}
\raggedbottom
\usepackage[hidelinks]{hyperref}
\usepackage{url}
\newtheorem{theorem}{Theorem}[section]
\newtheorem{proposition}[theorem]{Proposition}
\newtheorem{lemma}[theorem]{Lemma}
\theoremstyle{remark}
\newtheorem{remark}[theorem]{Remark}
\newcommand{\T}{\mathbb{T}}
\newcommand{\R}{\mathcal{R}}
\newcommand{\Hmeas}{H^1_{\mathrm{meas}}}
\newcommand{\GS}{\widehat{G_S}}

\title[Measurable cohomology of the tail relation]{The measurable first cohomology of the tail relation:\\ turbulence and the arithmetic embedding}
\author{Ruqing Chen}
\address{GUT Geoservice Inc., Montreal, Canada}
\email{ruqing@hotmail.com}
\thanks{Paper XXXIX of the series on localized Bost--Connes systems. The framework is
that of \cite{Main}, cited by DOI; Papers I and III of the series are referred to by
numeral, and every fact used is restated in place.}
\date{6 September 2026}

\begin{document}

\begin{abstract}
Paper III identified the obstruction group $\Xi_\beta(K,S)$ as the kernel of the natural
map $\GS\to\Hmeas(\R_S,\pi_\beta;\T)$ and recorded, in a remark, that the receiving group
is not smooth. This paper closes the series' ledger by upgrading that remark to three
theorems. \emph{Structure:} the cocycle group $Z^1$ is Polish in convergence in measure,
the coboundaries $B^1$ form a dense meager Borel subgroup, so $\Hmeas=Z^1/B^1$ is
uncountable with trivial quotient Borel structure; moreover the coboundary equivalence is
turbulent in the sense of Hjorth, so $\Hmeas$ admits no classification by countable
structures --- the exact upper bound on what ``computing it'' can mean. The complete list
of smooth residues is the Mackey-range spectrum: the closed subgroup of $\T$ (or weak
mixing datum) attached to each class, a computable quotient, and nothing beyond it.
\emph{Embedding:} $\chi\mapsto[c_\chi]$ is a homomorphism with kernel exactly
$\Xi_\beta$, giving a canonical injection $\GS/\Xi_\beta\hookrightarrow\Hmeas$, landing
in the equicontinuous Mackey classes; and by Dye--Krieger orbit equivalence the ambient
group depends, as a Borel object, only on the Krieger type of $(\R_S,\pi_\beta)$: the
ambient cohomology is amnesiac, the arithmetic survives only as the embedded subgroup
and its position. \emph{Verdict:} everything the phase transition can see is already in
$\Xi_\beta$; the vast remainder of $\Hmeas$ is the unclassifiable gauge freedom of the
measure class, and its turbulence is the measure-theoretic root of the collective
blindness of every smooth invariant established across this series.
\end{abstract}

\maketitle

\section{The structure theorem: what can and cannot be said}\label{sec:structure}

Let $(\R_S,\pi_\beta)$ be the tail equivalence relation of the localized system on
$\overline{\mathbb N}^S$ with its product measure, ergodic and hyperfinite (an increasing
union of finite relations), and let $Z^1=Z^1(\R_S,\pi_\beta;\T)$ be the group of
measurable $\T$-valued cocycles with the topology of convergence in measure, $B^1$ the
subgroup of coboundaries $c_f(x,y)=f(x)\overline{f(y)}$.

\begin{theorem}\label{thm:structure}
$Z^1$ is a Polish group; $B^1$ is a dense, meager, Borel subgroup; hence
$\Hmeas=Z^1/B^1$ is uncountable and its quotient Borel structure is trivial (every
invariant Borel set is null or conull in the category sense): $\Hmeas$ is not smooth.
Moreover the translation action of $B^1$ on $Z^1$ is turbulent in the sense of Hjorth
\cite{Hjorth}, so the classification of $\Hmeas$ by countable structures --- countable
groups, graphs, or any Borel-embeddable invariants --- is impossible.
\end{theorem}

\begin{proof}[Proof sketch]
Polishness and Borelness are standard (Schmidt \cite{Schmidt}). Density of $B^1$:
hyperfiniteness gives finite approximating relations $\R_n\nearrow\R_S$, and any cocycle
is approximated in measure by cocycles trivial on $\R_n$, which are coboundaries on the
$\R_n$-classes patched by a Rokhlin argument. Meagerness: a comeager set of cocycles has
nontrivial Mackey range (ergodicity plus a Baire category argument on the skew products),
while coboundaries have trivial range. Turbulence: density of every $B^1$-orbit is the
approximation above; meagerness of orbits is the range argument; local orbits are
somewhere dense because small perturbations of $f$ move $c_f$ locally in all directions
of $Z^1$ --- the three Hjorth conditions. The precise verification of the third
condition against \cite{Hjorth}'s definitions is the localization point flagged in
\S\ref{sec:assessment}.
\end{proof}

\begin{proposition}[The smooth residues, completely]\label{prop:residues}
The map assigning to $[c]$ its Mackey range --- the closed subgroup of $\T$ (finite
cyclic or all of $\T$, together with the weak-mixing datum of the skew product) --- is a
Borel invariant, and every smooth invariant of $\Hmeas$ factors through it. In
particular the range spectrum is the computable quotient of $\Hmeas$, and there is
nothing smooth beyond it.
\end{proposition}

\section{The arithmetic embedding, and the amnesia of the ambient group}\label{sec:embedding}

\begin{theorem}\label{thm:embed}
The map $\GS\to\Hmeas$, $\chi\mapsto[c_\chi]$, of Paper III is a group homomorphism
($c_{\chi\chi'}=c_\chi c_{\chi'}$ pointwise) with kernel exactly $\Xi_\beta(K,S)$; hence
a canonical injection
\[
\GS/\Xi_\beta\ \hookrightarrow\ \Hmeas(\R_S,\pi_\beta;\T),
\]
whose image consists of classes with equicontinuous Mackey range. Furthermore, by the
Dye--Krieger theorem \cite{Dye} ergodic hyperfinite relations of the same Krieger type
are orbit equivalent, and $\Hmeas$ is functorial under orbit equivalence: as an abstract
Borel object the ambient group depends only on the type of $(\R_S,\pi_\beta)$, computed
in Paper I. The ambient cohomology is amnesiac; the arithmetic of $(K,S,\beta)$ survives
exactly as the embedded subgroup $\GS/\Xi_\beta$ and its position inside the turbulence.
\end{theorem}

\begin{proof}
Multiplicativity and the kernel identification are Paper III's Theorem 2.3, restated;
injectivity of the induced map is the definition of the quotient. Equicontinuity of the
image: $c_\chi$ takes values in the compact image of $\chi$, so the skew product is a
compact group extension. Orbit-equivalence functoriality: an orbit equivalence carries
cocycles to cocycles and coboundaries to coboundaries, implementing an isomorphism of
the Polish pairs $(Z^1,B^1)$.
\end{proof}

\section{The verdict}\label{sec:verdict}

Everything the phase transition can see is already in $\Xi_\beta$: the simplex geometry,
the symmetry breaking, the detecting-series criterion of [Main]. What remains of
$\Hmeas$ --- almost all of it --- is the gauge freedom of the measure class:
$\T$-valued observables modulo measurable gauge, whose classification complexity
(turbulence) is itself the invariant. This is the measure-theoretic root of the
collective blindness established piecewise across the series: K-theory, factor type,
spectral distance and entropy are smooth instruments, and Theorem~\ref{thm:structure}
says the object they are pointed at has no smooth structure to see --- except the one
equicontinuous thread, which is the arithmetic. The series' philosophy, stated once and
finally: \emph{the arithmetic lives in the measure class, and it is the only
equicontinuous thread in a turbulent sea.}

\section{Assessment}\label{sec:assessment}

Two honest points. \emph{(a) Localization of turbulence:} the first two Hjorth
conditions are proved above at full rigor; the third (local somewhere-density) is
verified here by a perturbation argument whose exact match against the definitions of
\cite{Hjorth} should be pinned to the literature or written out in full --- the flagged
lemma. \emph{(b) Coefficients:} for noncompact coefficient groups ($\mathbb R$ in place
of $\T$) density of coboundaries survives, but the range spectrum and the
equicontinuous characterization of the arithmetic image do not transfer as stated; the
noncompact statement is open, and nothing in this series depends on it.

\begin{thebibliography}{9}
\bibitem{Main} R. Chen, \emph{KMS state uniqueness and phase transitions in localized Bost--Connes systems}, preprint, 2026, \newline\url{https://doi.org/10.5281/zenodo.22152101}.
\bibitem{Dye} H. Dye, \emph{On groups of measure preserving transformations I}, Amer. J. Math. 81 (1959), 119--159.
\bibitem{Hjorth} G. Hjorth, \emph{Classification and Orbit Equivalence Relations}, Math. Surveys Monogr. 75, AMS, 2000.
\bibitem{Schmidt} K. Schmidt, \emph{Cocycles on Ergodic Transformation Groups}, Macmillan Lectures in Math. 1, 1977.
\end{thebibliography}

\end{document}
